% Generated from accepted certificates; disjoint first successful tests.
\begin{table}[ht]
\centering
\begin{tabular}{p{.66\linewidth}rr}
\toprule
First successful top-level test, in dispatch order & First use $m$ & Rings\\
\midrule
Formal-codegree $d$-number failure & $1$ & $872\,625$\\
Rank-two moment obstruction & $2$ & $1\,273$\\
Non-self-dual quadratic moments & $2$ & $127$\\
Diagonal involution moments & $3$ & $126$\\
$C_3$ near-group root length & $9$ & $37$\\
Automatic near-group classification & $18$ & $3$\\
Cyclotomicity obstruction & $9$ & $84$\\
Fusion-subring dimension divisibility & $2$ & $28$\\
Residual-involution propagation & $5$ & $1$\\
Weak integrality and induction & $2$ & $1$\\
Two-dimension Galois reduction and induction & $3$ & $8\,126$\\
Signed sphericalization & $2$ & $7\,084$\\
\midrule
Total excluded & & $\ExcludedCount$\\
\bottomrule
\end{tabular}
\caption{Archived disjoint certificate allocation, displayed in the order
of the revised dispatcher. A model is recognized before these tests. First
use is the least exact multiplicity with a certificate of this type;
nested tests in signed records are not counted again. The sequence puts
high-yield direct obstructions before the expensive signed searches.}
\label{tab:proof-counts}
\end{table}
\noindent The finer ledger groups individual checks by their underlying
mathematical criterion. Its latest first use is at multiplicity 18.
In this fixed certificate allocation, all multiplicities from 19 to
$\MultiplicityBound$ require no additional criterion family. This describes
the completed run, not a proof that the collection of criteria is minimal.
For example, signed cyclotomicity is first decisive at multiplicity $95$,
but it applies the same cyclotomicity theorem used earlier for unsigned
characters. Signed-table parity is also used as an enumeration prefilter;
its removed candidate counts are recorded separately from the ring totals.
